\documentclass[12pt,a4paper]{article}
\usepackage{circuitikz}
\usepackage{amsmath}
\usetikzlibrary{circuits.logic.US,calc}


\begin{document}
\begin{center}
    \begin{circuitikz}[american,thick, x=1.3cm ]

    % \draw (0, 0) node[op amp] (op1) {};
    % \draw (5,-0.5) node[op amp] (op2) {};
    % \draw (op1.+) to[R , l=$3k\Omega$] (-4,-0.5) to[V , l=$10V$] (-4, -4);
    % \draw (-4,-4) -- ( 7,-4);
    % \draw (op1.-) -- (-1,0.5) -- (-1, 2) -- (1,2)-- ( 1,0);
    % \draw (-1.5,-0.5) to[R , l_=$10k\Omega$] (-1.5,-4); 
    % \draw (-1, -0.5) -- (-1,-2) to[R, l=$6k\Omega$] ( 3,-2) -- (3,2) to[R, l=$5k\Omega$]  (7,2)  --(7,-0.5) --(op2.out);
    %  \draw (7, -0.5) to[R, l=$5k\Omega$, i=$i_o$] (7,-4);
    % \draw (op1.out) to[R, l=$5k\Omega$] (3,0) -- (op2.-);
    % \draw (op2.+) -- (4, -1) -- (4,-4);
    % \draw (3,-4) node[ground]{};

    \draw (0,0) node[op amp] (op1){};
    \draw (op1.-) to[R ,l=$5k\Omega$] (-3,0.5) to[V , l=$2v$]  (-3,-3);
    \draw (op1.+) to (-1, -0.5) to[R, l=$10k\Omega$] (-1,-3) -- (-3,-3);
    \draw (-2,-3) node[ground]{};
    \draw (op1.-) -- (-0.9,1.5)  to[R, l=$20k\Omega$] (1.5,1.5)--(1.5,0);
    \draw (op1.out)  to[R, l=$40k\Omega$] (4,0);

    

    \draw (0,-8) node[op amp] (op2){};
    \draw (op2.-) to[R ,l=$10k\Omega$] (-3,-7.5) to[V , l=$3v$]  (-3,-13);
    \draw (op2.+) to (-1, -8.5) to[R, l=$30k\Omega$] (-1,-13) -- (-3,-13);
    \draw (-2,-13) node[ground]{};
    \draw (op2.out)  to[R, l=$80k\Omega$] (4,-8) to[R, l=$20k\Omega$] (4,-11) node[ground]{};
    \draw (-1, -10) to[R, l=$50k\Omega$] (1.5,-10)--(1.5,-8);


    \draw (5,-4) node[op amp] (op3){};
    \draw (op3.-) -- (4,0);
    \draw (op3.+) -- (4,-8);

    \draw (8.5,-4) node[op amp, rotate=180](op4){};

    \draw (op3.out)  to[R, l=$20k\Omega$, i=$i_o$] (op4.out);
    \draw (op4.-) to[R, l_=$20k\Omega$] (9.3,-7) to[V , l=$4v$] (9.3,-10) node[ground]{};


    \draw (4,0) to[R, l=$40k\Omega$] (6,0)-- (op3.out);
    \draw (op4.out) -- (7.5,-2)-- (9.5,-2)-- (op4.+);
\end{circuitikz}
\end{center}
\newpage




\begin{center}
\begin{circuitikz}
    % Set style to IEEE
    % \ctikzset{logic ports=us}
    
    % Define gates
    \node[and port] (AND1) at (0,0) {};
    \node[not port] (NOT1) at (0,-3) {};
    \node[not port] (NOT2) at (0,-7) {};
    \draw (AND1.in 1) -- ++(-4,0) node[left]{$P$};
    \draw (AND1.in 2) -- ++(-4,0) node[left]{$Q$};

    \draw (NOT2.in) -- ++ (-3,0) --++(0,7.3);

    \node[or port, anchor=in 2] (OR1) at (1,-6.5){};
    \draw (NOT2.out) -- (OR1.in 2);
    \draw (OR1.in 1) -- ++ (-4,0) --++ (0,5.7);
    \draw (NOT1.in) -- ++(-2.3,0);

    \node[and port] (AND2) at (3.5,-4){};
    \draw (AND2.in 1) -- (NOT1.out);
    \draw (AND2.in 2) --++ (-5.8,0);

    \node[or port] (OR2) at (6,-1){};
    \draw (AND1.out) -- (OR2.in 1);
    \draw (AND2.out) -- (OR2.in 2);

    \node[and port] (AND3) at (9, -2.5){};
    \draw (AND3.in 2) -- (OR1.out);
    \draw (AND3.in 1) -- (OR2.out);
    \draw (AND3.out) --++ (0.5,0);



    
    % Connect inputs and outputs
    % \draw (OR1.out) -- (AND1.in 1);
    % \draw (NOT1.out) -- (AND1.in 2);
    
    % % Label inputs/outputs
    % \draw (OR1.in 1) -- ++(-1,0) node[left]{$A$};
    % \draw (OR1.in 2) -- ++(-1,0) node[left]{$B$};
    % \draw (NOT1.in) -- ++(-1,0) node[left]{$C$};
    % \draw (AND1.out) -- ++(1,0) node[right]{$Out$};
\end{circuitikz}
\end{center}
\end{document}